\documentclass{report}
\usepackage{amsmath,amssymb}
\setlength{\parindent}{0mm}
\setlength{\parskip}{1em}
\begin{document}
\begin{center}
\rule{6in}{1pt} \
{\large Cao Lu \\
{\bf HyGA: A Hybrid Geometric+Algebraic Multigrid Solver for Weighted-Residual Methods with Hierarchical Meshes}}

Dept of Applied Math \& Stat  \\ Stony Brook University \\ Stony Brook \\ NY 11794
\\
{\tt clu@ams.sunysb.edu}\\
Xiangmin Jiao\\
Xiongfei Wei\end{center}

Multigrid methods are efficient and effective for solving large-scale
sparse linear systems, especially those from the discretizations of
partial differential equations (PDEs). Generally speaking, multigrid
may be classified into \emph{geometric} (\emph{GMG}) and \emph{algebraic}
(\emph{AMG}), each of which has advantages and disadvantages. The
primary advantages of GMG include its better convergence with smooth
solutions, better efficiency in terms of computational time, and also
better efficiency in storage especially with a matrix-free implementation.
However, the GMG tends to have difficulties for non-smooth problems
and irregular domains, and it is difficult to implement and to scale
with unstructured meshes, especially for domains with complex topologies.
The primary advantages of AMG include its simplicity and flexibility,
and its robustness for non-smooth solutions. However, AMG tends to
be more expensive, especially for its setup stage. It also tends to
generate denser matrices than GMG and cannot be easily implemented
in a matrix-free fashion, so it has higher memory and computational
costs. A more serious problem is that AMG has no guarantee on the
condition numbers of the matrices at the coarse level, which may lead
to non-convergence.

The goal of this work is to develop a novel and general \emph{hybrid
geometric+algebraic multigrid method}, or \emph{HyGA}, which will
combine the advantages and overcome the disadvantages of the GMG and
AMG. We focus on linear systems arising from a weighted residual
method for linear PDEs over unstructured meshes. Linear PDEs cover
a wide range of mathematical models, such as Poisson equations, and
the weighted residuals are general discretization techniques, which
can unify many commonly used finite elements and finite differences
by defining different weighting functions. In addition, unstructured
meshes are very general in resolving complex geometries. An important
contribution of this work is a new, unified derivation of restriction
and prolongation operators for weighted residuals with hierarchical
basis functions. This rigorous analysis coupled with the generality
of our formulation will allow us to address a large class of linear
systems arising from arising from scientific and engineering applications.

At the algorithmic level, the main contribution of this work is the
hybrid multigrid solver \emph{HyGA}, which combines a high-quality
hierarchical mesh generator, a geometric multigrid solver with a multilevel
weighted residual formulation, and an algebraic multigrid solver at
the coarsest levels. Our proposed algorithm may be summarized as follows.
Our hierarchical mesh generator starts from a good-quality coarse
unstructured mesh that is sufficiently accurate for representing the
topology and for reconstructing the geometry of the domain. It iteratively
refines the mesh with guaranteed mesh quality (by uniform refinements)
and geometric accuracy (by high-order boundary reconstruction). Because
this refinement involves only local operations, it is efficient and
scalable. We apply GMG with a multilevel weighted-residual formulation
on these hierarchical meshes using our unified prolongation and restriction
operators, which are coupled with PDE rediscretizations over coarse
meshes. This combination allows efficient matrix-free implementations
on the coarse levels, and also effective control of the sparsity and
the condition numbers of the resulting matrices. At the coarsest level,
we employ the classical AMG, which allow us to resolve complex topologies
easily and handle non-smooth solutions robustly. In addition, the
relatively high setup cost and memory requirement of AMG become negligible
due to the significantly reduced problem size. Therefore, our approach
effectively combines the rigor, high accuracy and runtime-and-memory
efficiency of GMG with the flexibility of AMG, and at the same time
it is relatively easy to implement.

We present numerical experiments to demonstrate the effectiveness
of HyGA compared with GMG and AMG for 2-D and 3-D problems. Our experimental
results show that both HyGA and GMG deliver better convergence and
better efficiency than AMG for weighted residual methods over unstructured
meshes. In addition, HyGA delivers similar, and sometimes better efficiency
and convergence rate than GMG due to the flexibility of AMG at the
coarse levels. In summary, our theoretical analysis and numerical
experimentations demonstrate that HyGA is an effective and robust
multigrid approach for weighted-residual methods with hierarchical
meshes.


\end{document}
